% CV — Maxence Pequeno
% Compile avec pdfLaTeX (Overleaf : compilateur par défaut).
\documentclass[10pt,a4paper]{article}

\usepackage[utf8]{inputenc}
\usepackage[T1]{fontenc}
\usepackage[french]{babel}
\usepackage[default]{sourcesanspro}
\usepackage[a4paper,top=1.3cm,bottom=1.3cm,left=1.6cm,right=1.6cm]{geometry}
\usepackage{xcolor}
\usepackage{titlesec}
\usepackage{enumitem}
\usepackage{tabularx}
\usepackage{fontawesome5}
\usepackage[hidelinks]{hyperref}

\definecolor{accent}{HTML}{1F3A5F}
\definecolor{muted}{HTML}{5A5F6B}

\pagestyle{empty}
\setlength{\parindent}{0pt}
\frenchsetup{StandardLists=true} % laisse enumitem gérer les puces

% --- Titres de section : petites capitales + filet ---------------------------
\titleformat{\section}
  {\color{accent}\large\bfseries\scshape}{}{0pt}{}[\vspace{-0.6ex}{\color{accent}\titlerule[0.6pt]}]
\titlespacing*{\section}{0pt}{1.6ex}{1.1ex}

\setlist[itemize]{leftmargin=1.1em, itemsep=0.15ex, topsep=0.4ex, label={\color{accent}\textbullet}}

% --- Une entrée : titre / dates, puis organisme / lieu ----------------------
% \entry{Poste}{Dates}{Organisme}{Lieu}
\newcommand{\entry}[4]{%
  \textbf{#1}\hfill{\color{muted}#2}\\
  {\color{accent}#3}\hfill{\color{muted}\small #4}\par\vspace{0.2ex}}

% \stack{technos}
\newcommand{\stack}[1]{{\color{muted}\small\textit{#1}}\par}

% \project{Nom}{Lien ou contexte}{Description}{Technos}
\newcommand{\project}[4]{%
  \textbf{#1} {\color{muted}\small — #2}\\
  #3\\
  \stack{#4}\vspace{0.7ex}}

\begin{document}

% ============================== EN-TÊTE ======================================
\begin{center}
  {\fontsize{24}{28}\selectfont\bfseries\color{accent} Maxence Pequeno}\\[0.8ex]
  {\large Développeur full-stack — Alternant en BUT Informatique}\\[1.2ex]
  \small
  \faMapMarker*\ Reims (51) \quad
  \faEnvelope\ \href{mailto:maxence.pequeno@icloud.com}{maxence.pequeno@icloud.com} \quad
  \faGlobe\ \href{https://pequenomaxence.fr}{pequenomaxence.fr} \\[0.5ex]
  \faGithub\ \href{https://github.com/MaxencePeq}{github.com/MaxencePeq} \quad
  \faLinkedin\ \href{https://www.linkedin.com/in/maxence-pequeno-671581329/}{linkedin.com/in/maxence-pequeno} \quad
  \faCar\ Permis B, véhiculé
\end{center}

% ============================== PROFIL =======================================
\section{Profil}
Étudiant en troisième année de BUT Informatique, en alternance chez Legras Industries.
Je développe des applications web complètes, du recueil du besoin à la mise en production :
back-end Symfony, FastAPI et Laravel, front-end React et Next.js.
Daltonien, je porte une attention particulière à l'accessibilité des interfaces.

% ============================== EXPÉRIENCE ===================================
\section{Expérience professionnelle}

\entry{Développeur full-stack — Alternance}{2026 -- 2027}{Legras Industries}{}
\begin{itemize}
  \item Interface entre plusieurs équipes métier et l'informatique : recueil des besoins,
        conception et livraison d'outils internes.
  \item Intranet pour les équipes : génération automatique de fichiers
        au format XML à partir du croisement de plusieurs sources de données, supprimant toute saisie manuelle.
\end{itemize}
\stack{Python, FastAPI, Jinja2, SQLite, Git}
\vspace{1ex}

\entry{Développeur web — Stage}{Avril -- Mai 2026}{Med'émérite}{Start-up e-learning}
\begin{itemize}
  \item Plateforme de cours en ligne (QCM, flashcards) : recueil des besoins,
        refactorisation et optimisation du code existant.
\end{itemize}
\stack{Laravel, React, Docker, MySQL}
\vspace{1ex}

\entry{Développeur web — Projet client}{2026}{Lebel Ouvrage — \href{https://lebelouvrage.bzh}{lebelouvrage.bzh}}{Arzal (56)}
\begin{itemize}
  \item Site vitrine d'un artisan installateur de portails, réalisé en binôme et en production depuis mai 2026.
  \item Formulaire de devis en plusieurs étapes (type de demande, matériel, commune via l'API
        geo.api.gouv.fr, photos) envoyant automatiquement un dossier complet par mail à l'artisan.
\end{itemize}
\stack{Symfony, Twig, JavaScript, Tailwind CSS, Git}

% ============================== PROJETS ======================================
\section{Projets}

\project{Homelab auto-hébergé}{personnel, 2026}
  {Rack de Raspberry Pi monté à la main : DNS Pi-hole, applications conteneurisées sous Docker,
   reverse proxy et HTTPS avec Caddy, exposition via Cloudflare Tunnel sans ouverture de port.
   Héberge mon portfolio.}
  {Docker, Caddy, Cloudflare, Pi-hole, Linux, SSH}

\project{Portfolio}{personnel, \href{https://github.com/MaxencePeq/portfolio}{GitHub}}
  {Migration de React vers Next.js, sans dépendance hors du framework ; poids de la page
   d'accueil réduit de 13\,Mo à moins de 400\,Ko. Image Docker multi-stage.}
  {Next.js, TypeScript, Tailwind CSS, Docker}

% ============================== AUTO-FORMATION ===============================
\section{Auto-formation}
\begin{itemize}
  \item \textbf{Go et Rust} — apprentissage en autodidacte : typage statique, concurrence
        (goroutines, channels), modèle d'ownership et d'emprunt de Rust.
  \item \textbf{Programmation bas niveau} — pointeurs, allocation et gestion manuelle de la mémoire,
        pour comprendre ce qui se passe sous les langages de haut niveau.
\end{itemize}

% ============================== FORMATION ====================================
\section{Formation}

\entry{BUT Informatique — 3\textsuperscript{e} année en alternance}{2024 -- 2027}
      {Université de Reims Champagne-Ardenne (URCA)}{Reims}
\vspace{0.6ex}
\textbf{Baccalauréat général — spécialités Science Numérique et Mathématiques}\hfill{\color{muted}2024}\par

% ============================== COMPÉTENCES ==================================
\section{Compétences}
\begin{tabularx}{\linewidth}{@{}>{\bfseries}l X@{}}
  Langages       & PHP, Python, TypeScript, JavaScript, SQL, C\# \\
  Back-end       & Symfony, API Platform, Laravel, FastAPI, Jinja2 \\
  Front-end      & React, Next.js, Tailwind CSS, Bootstrap \\
  Données, infra & MySQL, SQLite, Docker, Linux, Git, Caddy \\
  Méthodes       & Merise (MCD/MPD), UML, Scrum, cahier des charges, diagramme de Gantt \\
  Outils         & JetBrains, Neovim, Figma, Claude Code, MCP \\
\end{tabularx}

% ============================== DIVERS =======================================
\section{Langues et centres d'intérêt}
\begin{tabularx}{\linewidth}{@{}>{\bfseries}l X@{}}
  Langues   & Français (langue maternelle), anglais (C1) \\
  Intérêts  & Piano (depuis 2021), littérature, peinture aquarelle \\
\end{tabularx}

\end{document}
